\par\noindent\SourceLeaf{50}%
\hspace*{\parindent}Cum sit igitur $\cos\varphi^{2m-1}$ series, cuius terminus initialis
\[
\frac{1}{2^{2m-2}}\,{}^{2m-1}\mathfrak{B}^{(m)}\cos\varphi
\]
erit, substituendo $90^0-\varphi$ in locum ipsius $\varphi$, $\sin\varphi^{2m-1}$
terminus initialis
\[
\frac{1}{2^{2m-2}}\,{}^{2m-1}\mathfrak{B}^{(m)}\sin\varphi
\]
et generatim, eadem substitutione $\sin\varphi^{2m-1}$ terminus $r^{\text{tus}}$
\[
= \frac{1}{2^{2m-2}}\,{}^{2m-1}\mathfrak{B}^{(m+r)}\cos(2r+1)(90^0-\varphi).
\]
Est autem
\[
\begin{aligned}
\cos(2r+1)(90^0-\varphi) &= \cos\big((2r+1)90^0 - (2r+1)\varphi\big)\\
&= \cos(2r+1)90^0 . \cos(2r+1)\varphi\\&\quad + \sin(2r+1)90^0 . \sin(2r+1)\varphi.
\end{aligned}
\]
Iam vero, cum sit
\[
\sin(2r+1)90^0 = (-1)^r, \qquad \cos(2r+1)90^0 = 0,
\]
fiet haec quantitas
\[
= (-1)^r\sin(2r+1)\varphi,
\]
et terminus ipsius $\sin\varphi^{2m-1}$ generalis $r^{\text{tus}}$
\[
= \frac{(-1)^r}{2^{2m-2}}\,{}^{2m-1}\mathfrak{B}^{(m+r)}\sin(2r+1)\varphi.
\]

Si vero in serie $\cos\varphi^{2m}$ exhibente in locum ipsius $\varphi$
substituas $90^0-\varphi$, terminus initialis non afficitur, generalis
$r^{\text{tus}}$, qui est
\[
\frac{1}{2^{2m-1}}\,{}^{2m}\mathfrak{B}^{(m+r)}\cos 2r\varphi,
\]
mutatur ita, ut
\[
\begin{aligned}
\cos 2r(90^0-\varphi) &= \cos(2r\,90^0 - 2r\varphi)\\
&= \cos 2r\,90^0 . \cos 2r\varphi + \sin 2r\,90^0 . \sin 2r\varphi
\end{aligned}
\]
contineat. Est autem
\[
\sin 2r\,90^0 = 0, \qquad \cos 2r\,90^0 = (-1)^r,
\]
qua re mutatur ista expressio in $(-1)^r\cos 2r\varphi$, fitque ipsius
$\sin\varphi^{2m}$ terminus $r^{\text{tus}}$
\[
\frac{(-1)^r}{2^{2m-1}}\,{}^{2m}\mathfrak{B}^{(m+r)}\cos 2r\varphi.
\]
